\newpage
\section{Tujuan Magang Industri}

Pelaksanaan Magang Industri tidak hanya bertujuan untuk memenuhi kewajiban
akademik, tetapi juga menjadi sarana bagi mahasiswa untuk mempersiapkan diri
menghadapi dunia kerja profesional. Secara khusus, tujuan magang industri ini
adalah:

\begin{enumerate}[noitemsep]
      \item Menerapkan dan mengembangkan kompetensi teknis dalam pengembangan
            \textit{software} pada lingkungan profesional yang mengadopsi standar industri,
            mencakup alur \textit{software development}, pengelolaan \textit{version
                  control}, penerapan \textit{clean code}, serta pemecahan masalah teknis secara
            mandiri maupun kolaboratif dalam tim.

      \item Memahami dan mengikuti proses \textit{Software Development Life Cycle} (SDLC)
            secara langsung, mulai dari penerimaan task, implementasi, \textit{code
                  review}, pengujian, hingga proses \textit{deployment} ke lingkungan
            \textit{staging} dan \textit{production}.

      \item Mengembangkan kemampuan teknis dalam pengembangan aplikasi menggunakan
            framework Flutter dan React (\textit{Next.js} dan TypeScript), serta memahami
            proses integrasi layanan \textit{back-end} berbasis Django/Python dalam konteks
            sistem manajemen klinik (\textit{healthcare information system}) berskala
            \textit{enterprise}.

      \item Berperan aktif sebagai \textit{Full-Stack Developer} dalam tim lintas fungsi,
            termasuk melakukan koordinasi dengan tim \textit{Back-End}, \textit{Quality
                  Assurance} (QA), \textit{Project Manager} (PM), dan anggota tim lainnya untuk
            menghasilkan fitur yang sesuai dengan \textit{requirement} serta standar
            kualitas perusahaan.

      \item Memperoleh pemahaman nyata mengenai dinamika kerja profesional di industri
            teknologi informasi, termasuk budaya kerja, sistem evaluasi kinerja berbasis
            \textit{Key Performance Indicator} (KPI), penggunaan alat kolaborasi seperti
            OpenProject dan GitLab, serta praktik pengembangan produk digital secara
            \textit{end-to-end}.

      \item Mengembangkan kemampuan adaptasi terhadap perubahan teknologi dan kebutuhan
            bisnis melalui pengalaman transisi \textit{tech stack} dari Flutter ke React
            serta keterlibatan dalam pengembangan layanan \textit{back-end}, sehingga
            diperoleh pemahaman yang lebih komprehensif mengenai proses pengembangan sistem
            modern.
\end{enumerate}